% concatenated from main.tex
\documentclass[11pt, reqno]{amsart}
\usepackage[english]{babel}
\usepackage[table]{xcolor}
\usepackage[utf8]{inputenc}
\usepackage{amsmath, amsfonts, amsthm, amscd, amssymb, fullpage, upgreek, enumerate, comment, siunitx, enumitem, graphicx, mathrsfs, mathtools}
\usepackage{thmtools}
\usepackage{thm-restate}
\usepackage{relsize}
\usepackage{caption}
\captionsetup{width=0.45\textwidth}

\usepackage[parfill]{parskip}
\usepackage{hyperref}

\newtheorem{theorem}{Theorem}
\newtheorem{proposition}[theorem]{Proposition}
\newtheorem{corollary}[theorem]{Corollary}
\newtheorem{lemma}[theorem]{Lemma}
\newtheorem{definition}[theorem]{Definition}
\newtheorem{step}{Step}
\newtheorem{claim}{Claim}

\newtheorem{problem}{Problem}

\newtheorem*{remark}{Remark}
\newtheorem*{notation}{Notation}

\numberwithin{equation}{section}

\DeclareMathOperator{\sym}{sym}
\DeclareMathOperator{\ad}{ad}
\DeclareMathOperator{\GL}{GL}
\DeclareMathOperator{\res}{res}

\newcommand{\bburl}[1]{\textcolor{blue}{\url{#1}}}

\newcommand{\blue}[1]{\textcolor{blue}{(\bf{#1})}}

\newcommand{\emaillink}[1]{\textcolor{blue}{\href{mailto:#1}{#1}}}

\renewcommand{\baselinestretch}{1}
\newcommand{\murl}[1]{\href{mailto:#1}{\textcolor{blue}{#1}}}
\newcommand{\hr}[1]{\href{#1}{\url{#1}}}
\newcommand{\fix}[1]{{\color{red} Fix: [#1]}}
\newcommand{\loweq}{\mathrel{\underset{\sim}
{\leqslant}}}\newcommand{\textcolour}[2]{\textcolor{#1}{#2}}
\newcommand{\cube}{\{0, 1, 2\}^n}

\newcommand\modd{\operatorname{mod}}
\newcommand\new{\operatorname{new}}

\newcommand{\nocontentsline}[3]{}

\renewcommand{\footskip}{25pt}

\setcounter{MaxMatrixCols}{12}

\title{A Generalisation of Sperner's Theorem \\ Using Weighted Chains}
\author{Ya\"el Dillies}
\thanks{Department of Mathematics, Stockholm University. Email: yael.dillies@math.su.se}

\author{Matthew Johnson}
\thanks{Independent researcher, London, UK. Email: mjohnson31415926@gmail.com}

\author{Aleksandra Kowalska}
\thanks{Mathematical Institute, University of Oxford. Email: aleksandra.kowalska@seh.ox.ac.uk}

\begin{document}

\begin{abstract}
    We find the (unique) largest subset of $\{0, 1, 2\}^n$ such that it contains no two elements, one of which is coordinatewise greater than the other, but strictly greater on at most $k$ coordinates. To do so, we decompose the cube into weighted chains.

    In Appendix \ref{appendix_sperner} we present a new proof of Sperner's theorem we found while working on this problem.
\end{abstract}

\maketitle

\section{The problem statement and initial motivation}

In this paper, we consider the following problem:

\begin{problem}
\label{the_main_problem}
    For $d, n, k \geqslant 1$, what is the largest subset $A \subseteq \{0, 1, \dots, d\}^n$ such that there are no two distinct elements $x, y \in A$ with $x_i \leqslant y_i$ for all $1 \leqslant i \leqslant n$ and $x_i < y_i$ for at most $k$ coordinates $i$?
\end{problem}

Hence, $A$ is allowed to have two elements, one of which is coordinatewise greater than the other, but only if the larger one is strictly greater on more than $k$ coordinates.

We note that for $k=1$, the problem is very simple: $|A| \leqslant (d+1)^{n-1}$, and for any $a$ the set
\begin{equation}
    \Big\{ x: |x| \equiv a \ (\modd d+1) \Big\}
\end{equation}
has this size (where $|x|=\sum_i x_i$).

On the other hand, for $\frac{n}{2} \leqslant k \leqslant n$, the problem is also relatively simple: the largest sets are of the same size as
\begin{equation}
    B = \Big\{ x: |x| = \Big\lfloor \frac{dn}{2} \Big\rfloor \Big\},
\end{equation}
This also follows from \cite[Theorem 5.1]{tour_m_part_p_sperner_families}. Below we provide a proof at which we arrived independently (we note that the same proof would also work for cuboids $\{0, 1, \dots, d_1\} \times \dots \times \{0, 1, \dots, d_n\}$).

Let $X=\{0, \dots, d\}^{\lfloor \frac{n}{2} \rfloor}, Y=\{0, \dots, d\}^{\lceil \frac{n}{2} \rceil}$ and let us consider the decomposition $\{0, \dots, d\}^n=X \times Y$. For two tuples $(x, y), (x', y') \in \{0, \dots, d\}^n$ (where $x, x' \in X,$ $y, y' \in Y$), if $x=x'$ and $y \leqslant y'$ or $y=y'$ and $x \leqslant x'$, both cannot belong to $A$. Let us decompose $X$ and $Y$ into symmetric chains (as in the proof by de Bruijn, Tengbergen and Kruyswijk \cite{symmetric_chain_decomposition}) respectively $C_1, \dots, C_p$ and $D_1, \dots, D_q$. We cannot have more than $\min(|C_i|, |D_j|)$ elements $(x, y) \in A$ with $x \in C_i, y \in D_j$. We note that $B$ has exactly this number of such elements. Hence, we have
\begin{equation}
    |A|= \sum_{i, j} |A \cap (C_i \times D_j)| \leqslant \sum_{i, j} \min(|C_i|, |D_j|) =  \sum_{i, j} |B \cap (C_i \times D_j)| = |B|,
\end{equation}
so $B$ is an optimal such set.

However, the problem for the intermediate values of $k$, $2 \leqslant k < \frac{n}{2}$ is much harder. In this paper, we present the solutions for $d=1$ and $d=2$. In particular, the solution for $d=2$ (stated in the following theorem) is the main contribution of this paper.

\begin{theorem}
\label{the_main_theorem}
    For any $0 \leqslant k \leqslant n$, the unique largest subset $A$ of $\cube$ s.t. there are no two distinct elements $x, y \in A$ with $x \leqslant y$ coordinatewise and $x_i < y_i$ on at most $k$ coordinates, is
    \begin{equation*}
        \Big\{ x \in \cube: |x| \equiv n \ (\modd 2k+1) \Big\},
    \end{equation*}
    where $|x|=\sum_i x_i$.
\end{theorem}

In the next section, we describe related problems and research (and explain, why we use weighted chain decomposition, instead of just a chain decomposition). In Section \ref{section_d1}, we provide the proof in case $d=1$, which is an easier special case. It easily follows from Katona's \cite{katona} and from Grigg's \cite{griggs} (our proof was arrived at independently from these). Although it is not a novel result, we present our new proof for $d=1$ with weighted chains, as it provides a good motivation for the proof of the $d=2$ case (Theorem \ref{the_main_theorem}), which is presented in Section \ref{section_d2_overview}. The last section concludes the paper and discusses possible generalisations for $d>2$.

\subsection{Acknowledgements}
We would like to thank our supervisors Ben Green and Tom Bloom for very helpful discussions and support.

\section{Background and related research}

\subsection{Motivation for the problem}

Our initial motivation for studying the problem was trying to prove the following:

\begin{problem}
    Let $A_1, \dots, A_n$ be disjoint sets. What is the size of the largest family of subsets $\mathcal{F} \subseteq \mathcal{P}(A_1 \cup \dots \cup A_n)$ such that for any two $X, Y \in \mathcal{F}$ if $X \nsubseteq Y$ then we must have $X \cap A_i \neq Y \cap A_i$ for all $1 \leqslant i \leqslant n$?
\end{problem}

Our method of solving the general problem (using Symmetric Chain Decomposition) required the case $k=n$ of the problem presented in this paper. Having solved the special case with $k=n$, we tried to generalise it to other values of $k$, which motivated the work whose results are presented in this note.

The above problem has been solved in a paper \cite[Theorem 5.1]{tour_m_part_p_sperner_families} on various similar generalisations of Sperner's Theorem.

\subsection{A similar problem in a different metric}

Let us consider the following problem, which might sound similar to Problem \ref{the_main_problem}:

\begin{problem}
\label{another_problem}
    For any $k, n$ what is the largest subset $A \subseteq \{0, 1, \dots, d\}^n$ s.t. for any distinct $x, y \in A$ with $x_i \leqslant y_i$ for all $1 \leqslant i \leqslant n$ we must have $\sum_i |x_i - y_i| > k$?
\end{problem}

We note that both this problem and Problem \ref{the_main_problem} can be formulated in a very similar way, asking us to find the largest subset $A \subseteq \{0, 1, \dots, d\}^n$ s.t. for any distinct $x, y \in A$ with $x \leqslant y$ coordinatewise we must have $|x-y|>k$, where in the version above $|x|=\sum_i |x_i|$ and in the original version $|x|=|\{i : x_i \neq 0\}|$. Moreover, the cases with $d=1$ of both problems are identical. However, as we will soon discuss, these two problems are solved with rather different methods.

In \cite{katona}, Katona proved the following theorem, which solves the $d=1$ case of both problems:

\begin{theorem}
\label{katona_theorem}
    The largest possible size of a family of subsets $\mathcal{A} \subseteq 2^{[n]}$ s.t. there are no distinct $X, Y \in \mathcal{A}$ with $X \subseteq Y$ and $|Y \setminus X| \leq k$ is $\sum_{i \equiv \lfloor \frac{n}{2} \rfloor} \binom{n}{i}$.
\end{theorem}

The main part of Katona's proof is the following more general lemma (\cite[Theorem 1]{katona}):

\begin{lemma}
\label{katona_lemma}
    For a cuboid $\{0, 1, \dots, d_1\} \times \dots \times \{0, 1, \dots, d_n\}$ let us call a \emph{chain} a tuple of points $x_1 \leqslant x_2 \leqslant \dots \leqslant x_l$ and $\sum_j |x_{i+1, j} - x_{i, j}|=1$ for all $1 \leqslant i \leqslant l-1$ (where $x_{i,j}$ denotes the $j^{\text{th}}$ coordinate of point $x_i$ and the inequalities are coordinatewise).

    Then for any $l, d_1, \dots, d_n$ there exists some $m$ s.t. the cuboid $\{0, 1, \dots, d_1\} \times \dots \times \{0, 1, \dots, d_n\}$ can be decomposed into chains, each of which is either of length $l$ or shorter than $l$ and symmetrical around level $m$ (i.e., the sum of coordinates of its first and last element is $2m$).
\end{lemma}

Katona proves this lemma using a recursive construction, somewhat reminiscent of de Bruijn's, Tengbergen's and Kruyswijk's proof of the existence of Symmetric Chain Decomposition \cite{symmetric_chain_decomposition}.

We note that in particular, the lemma solves Problem \ref{another_problem}. Indeed, let
\begin{equation*}
    B_p := \Big\{ x \in \{0, 1, \dots, d_1\} \times \dots \times \{0, 1, \dots, d_n\} : |x| \equiv p \ (\modd k+1) \Big\}
\end{equation*}
(we note that for any $p$, $B_p$ satisfies the conditions of Problem \ref{another_problem}). Let us consider a decomposition of a cube into chains with $l=k+1$.

Since each chain passes through exactly one element of $B_m$, the number of chains is at most $|B_m|$.
On the other hand, since no chain can pass through more than one element of a set $A$, satisfying the problem's conditions, so $|A| \leqslant |B_m|$. Hence, $B_m$ is one of optimal sets, so $|B_m|=\max_p |B_p|$, which implies that $m=\lfloor \frac{dn}{2} \rfloor$ or $m=\lfloor \frac{dn}{2} \rfloor$ (this is a simple result, presented in Lemma \ref{layers_modulo_largest}). Hence, $B_{\lfloor \frac{dn}{2} \rfloor}$ is the optimal such set.

\subsection{Why the methods used to solve Problem \ref{the_main_problem} are different from those used for Problem \ref{another_problem}.}

Despite Problems \ref{the_main_problem} and \ref{another_problem} appearing similar, they require rather different approaches. In particular, none of methods for proving Problem \ref{another_problem} that we found seemed adaptable to Problem \ref{the_main_problem}, which was suggesting that to solve it, one needs to study the cube from a different angle.

The first method of solving Problem \ref{another_problem} is Katona's decomposition of the cube into chains that are either of length $k+1$ or symmetric around some $m$, as discussed above. If it is possible to construct a similar decomposition, but with chains of unlimited length and width at most $k$ (where a \emph{width} of a chain is the number of coordinates on which its first and last elements differ), then it could be used to solve Problem \ref{the_main_problem} analogously to how Problem \ref{another_problem} was solved above.

However, Katona's recursive (for $n$, the number of dimensions of the cube, increasing) construction of bounded-length chain decomposition does not seem to adapt to bounded-width chain decomposition. The main difference seems to lie in the fact that for the purposes of bounded-length decomposition, only the length of a chain matters (and not its shape, or which points specifically it passes through) - so knowing the length of a chain and where it ends is enough to be able to extend it without breaking the length condition. This is, however, not the case for a bounded-width chain decomposition, where to extend a chain while preserving its width condition, we require information on the shape of the chain (i.e., on which coordinates it has changed).

We also note a small `aesthetical' advantage of weighted chain decomposition over a standard chain decomposition: while the latter one is defined recursively, so how it looks like for larger dimensions might not be clear, the chains that appear in the former with positive weight are explicitly defined (only with weights defined recursively). Moreover, unlike the latter, the former decomposition is preserved under permuting coordinates.

Another method of solving Problem \ref{another_problem} is by using a method presented by D. Nagy and K. Nagy in \cite{nagys} (where Katona's Theorem \ref{katona_theorem} is one of results the authors use as an example for application of their method).

The reason the Nagys' method fails for Problem \ref{the_main_problem} is that adapting it would require that for any maximal set $A \subseteq \{0, 1, \dots, d\}^n$ satisfying our condition (from Problem \ref{the_main_problem}) and any ordered subset $(C_1, \dots, C_l) \subseteq \{0, 1, \dots, d\}^n$ s.t. $|C_{i+1}-C_i|=1$ of length $1+dn$, the set $A \cap C$ is the largest possible subset of $C$ satisfying our condition, which is not the case.

\section{The \texorpdfstring{$d=1$}{d=1} case}
\label{section_d1}

In this section, for any $1 \leqslant k \leqslant n$ we find the largest $A \subseteq \{0, 1\}^n$ such that $(A-A) \cap X = \emptyset$, where $X=\Big\{x \in \{0, 1\}^n: |x| \leqslant k\Big\}$. For $n$ even, the unique such set is
\begin{equation}
\label{b_definition}
    B = \Big\{ x \in \{0, 1\}^n: |x| \equiv \frac{n}{2} \ (\modd k+1) \Big\},
\end{equation}
and for $n$ odd the only two such sets are
\begin{equation}
    B_1 = \Big\{ x \in \{0, 1\}^n: |x| \equiv \Big\lfloor \frac{n}{2} \Big\rfloor \ (\modd k+1) \Big\} \text{  and  } B_2 = \Big\{ x \in \{0, 1\}^n: |x| \equiv \Big\lceil \frac{n}{2} \Big\rceil \ (\modd k+1) \Big\}.
\end{equation}
We note that the case $k=n$ immediately implies Sperner's theorem. This special case has the same proof idea, but slightly different (simpler) details of the proof and can be found in Appendix \ref{appendix_sperner}. To our knowledge, this is a new proof of Sperner's Theorem, not found in the literature.

\subsection{Definitions and notation}

First, let us introduce a couple of definitions that will make dealing with elements of the cube $\{0, 1\}^n$ easier.

\begin{definition}
    Let a \emph{layer} $L_i$ be the subset of the cube $\{0, 1\}^n$ of all elements with a specific number of $1$s, i.e.
    \begin{equation*}
        L_m=\Big\{ x \in \{0, 1\}^n: |x|=m \Big\}.
    \end{equation*}
    We will call the layers $L_m$ with $m \leqslant \frac{n}{2}$ \emph{lower layers,} and the layers with $m \geqslant \frac{n}{2}$ \emph{upper layers.} If $n$ if even, the middle layer can be called either.

    For $n$ even the layer $L_{\frac{n}{2}}$ and for $n$ odd the layers $L_{\lfloor \frac{n}{2} \rfloor}, L_{\lceil \frac{n}{2} \rceil}$ are called the \emph{middle layers.}
\end{definition}

We note that $|L_m|=\binom{n}{m}$. We also note that types are orbits under the group action of permuting coordinates.

\begin{definition}
    For $x, y \in \{0, 1\}^n$ we write $x \prec y$ if there exists a coordinate $j$ such that $x_i=y_i$ for all $i \neq j$ and $x_j=0$, $y_j=1$. We note that this relation isn't transitive (we don't need to make $\{0, 1\}^n$ into a poset structure for this proof).

    We call a \emph{chain} an ordered tuple $(x_1, x_2, \dots, x_l) \subseteq \{0, 1\}^n$ such that $x_1 \prec x_2 \prec \dots \prec x_l$.

    The \emph{length} of a chain is the number of its elements.
\end{definition}

We note that a tuple consisting of a single element is a chain. We also note that if $\gamma=(x_1, \dots, x_l)$ is a chain and $x_1 \in L_m$, then $x_2 \in L_{m+1}, \dots, x_l \in L_{m+l-1}$.

\begin{definition}
    We call a chain \emph{symmetric} if $\gamma=(x_1, \dots, x_l)$ and $x_1$ lies in layer $L_m$ (so $x_l$ lies in layer $L_{m+l-1}$) and $\frac{m+(m+l-1)}{2}=\frac{n}{2}$.
\end{definition}

Let us note that for $n$ even, all symmetric chains are of odd length, with the first element in layer $L_{\frac{n}{2}-i}$ and the last element in layer $L_{\frac{n}{2}+i}$ for some $i$. For $n$ odd, all symmetric chains are even, with the first element in layer $L_{\lfloor \frac{n}{2} \rfloor -i}$ and the last element in layer $L_{\lceil \frac{n}{2} \rceil + i}$.

\subsection{Weighted chains}

Let
\begin{multline}
    \mathcal{C} = \Big\{ \gamma: \gamma \text{ is a chain in } \{0, 1\}^n, |\gamma| = k+1\Big\}
    \\ \cup \Big\{ \gamma: \gamma \text{ is a symmetric chain in } \{0, 1\}^n, |\gamma| < k+1 \Big\}
\end{multline}
(where $|\gamma|$ denotes the number of elements in a chain $\gamma$). We note that each of the chains in $\mathcal{C}$ contains exactly one element of $B$.

The figure below shows two examples of $\mathcal{C}$. The dots on the drawings represent layers, and a brace between two dots means that all chains starting at one of these layers and finishing at the other one are in $\mathcal{C}$. The black dots represent the middle layer(s).

\begin{figure}[h!]
    \centering
    \begin{minipage}{0.45\textwidth}
        \centering
        \includegraphics[scale=0.7]{n10_k4.png}
        \caption{$n=10, k=3$}
        \label{d1_drawing1}
    \end{minipage}\hfill
    \begin{minipage}{0.45\textwidth}
        \centering
        \includegraphics[scale=0.55]{n13_k6.png}
        \caption{$n=13, k=5$}
        \label{d1_drawing2}
    \end{minipage}
\end{figure}

We note that if $n$ is even, $\mathcal{C}$ includes `singleton' chains consisting of elements of the middle layer, and for $n$ odd it includes chains of length two, consisting of all pairs $x \prec y$ for $x \in L_{\lfloor \frac{n}{2} \rfloor}, y \in L_{\lceil \frac{n}{2} \rceil}$. The lemma below is the crucial step of the proof for $d=1$.

\begin{lemma}
    If there exist strictly positive weights $w_{\gamma}$ for each $\gamma \in \mathcal{C}$ such that for any $x \in \{0, 1\}^n$, we have
    \begin{equation}
        w_x := \sum_{\substack{\gamma \in \mathcal{C}: \\ x \in \gamma}} w_{\gamma} = 1
    \end{equation}
    (we will call $w_x$ the \emph{induced} weight of $x$), then $B$ (for $n$ even) or $B_1$ and $B_2$ (for $n$ odd) are the only largest sets satisfying our condition.
\end{lemma}

First we will prove the lemma and then show that such an assignment of weights exists.

\begin{proof}
        Assuming that such an assignment of weights exists, we can easily finish the proof. First, let us do so for $n$ even. Indeed, let us assume that $A$ satisfies the conditions of the problem. Then $A$ has at most one element in any chain in $\mathcal{C}$ with positive weight, and $B$ (defined in Equation \ref{b_definition}) has exactly one element in each such chain. Hence, we have
    \begin{equation}
        |A| = \sum_{a \in A} 1 = \sum_{a \in A} w_a = \sum_{a \in A} \sum_{\substack{\gamma \in \mathcal{C}: \\ a \in \gamma}} w_{\gamma} = \sum_{\substack{\gamma \in \mathcal{C}: \\ A \cap \gamma \neq \emptyset}} w_{\gamma} \leqslant \sum_{\substack{\gamma \in \mathcal{C}}} w_{\gamma} = \sum_{\substack{\gamma \in \mathcal{C}: \\ B \cap \gamma \neq \emptyset}} w_{\gamma} = \sum_{b \in B} \sum_{\substack{\gamma \in \mathcal{C}: \\ b \in \gamma}} w_{\gamma} = \sum_{b \in B} w_b = |B|,
    \end{equation}
    so $B$ is one of the maximal sets satisfying the condition. Let us suppose that we have $|B|=|B'|$ for some other $B'$ satisfying the conditions of the problem. Since the weights of all chains are positive, $B'$ must have at least one element in all chains of $\mathcal{C}$. In particular, it needs to include the entire middle layer (as the singleton chains with elements of the middle layer are in $\mathcal{C}$).
    
    Now let us note that, for any chain of length $k+1$ with the first element in layer $L_{\frac{n}{2}-k-1}$ and the last element in layer $L_{\frac{n}{2}-1}$, at least one of its elements must be in $B'$. However, having any element from $L_{\frac{n}{2}-k}, \dots, L_{\frac{n}{2}-2}$ or $L_{\frac{n}{2}-1}$ in $B'$ would violate the condition, so the first element of the chain must be in $B'$. Reasoning this way for all such chains, we find that $L_{\frac{n}{2}-k-1} \subseteq B'$.
    
    Repeating this reasoning for further layers and on the other side of the cube, we find that $L_{\frac{n}{2}+i(k+1)} \subseteq B'$ for any $i$, and so $B \subseteq B'$, so in fact $B'=B$.
    
    For $n$ odd, the proof is completely analogous, except that at the beginning of showing the uniqueness, we need to show that either $L_{\lfloor \frac{n}{2} \rfloor} \subseteq B'$ or $L_{\lceil \frac{n}{2} \rceil} \subseteq B'$.
    
    To do this, let us assume without loss of generality that $x \in B'$ for some $x \in L_{\lfloor \frac{n}{2} \rfloor}$, and we will show that $L_{\lfloor \frac{n}{2} \rfloor} \subseteq B'$. Indeed, for any $y \in L_{\lfloor \frac{n}{2} \rfloor}$ which differs from $x$ in exactly two coordinates, there exists $z \in L_{\lceil \frac{n}{2} \rceil}$ such that $x \prec z, y \prec z$, so $z \notin B'$, but also one of $y, z$ must be in $B'$, so $y \in B'$. Hence, since we can get from any element of $L_{\lfloor \frac{n}{2} \rfloor}$ to any other by changing two coordinates at a time (by swapping a $0$ with a $1$ a few times), we in fact have $L_{\lfloor \frac{n}{2} \rfloor} \subseteq B'$. The rest of the proof is analogous to the case of $n$ even.
\end{proof}

Now, we only need to show that an assignment of positive weights to chains in $\mathcal{C}$, such that the induced weight of any element of the cube is $1$, is possible. We do so in the remainder of this section, first defining the weights and then showing that they are positive.

\subsection{Assigning chain weights so that the induced weights are all \texorpdfstring{$1$}{1}}

We will say that a chain whose element furthest from the middle of the cube (which must be either its first or last element) lies in layer $L_m$ \emph{starts} at that layer (the symmetric chains can be said to start at either of their ends).

We will assign the same weight to all chains in $\mathcal{C}$ starting at the same layer. Let $W(m)$ be the total weight of all chains starting at layer $L_m$ (where $W(m)$ is the common weight of these chains multiplied by the number of such chains).

We can assign the weights recursively, from chains starting the furthest from the middle, to those starting the closest, ensuring at each step that all points have induced weight $1$, as follows:

Let $W(0)=W(n)=1$ (as there is only one point in both of these layers). Then, let
\begin{equation}
    W(m)=\binom{n}{m}-\sum_{\substack{i: \ |i-m| \leqslant k-1, \\
    |i-\frac{n}{2}|>|m-\frac{n}{2}|}} W(i)
\end{equation}
for $m=2, n-2, 3, n-3, \dots$. We note that at the moment of assigning $W(m)$, the weights of all other chains passing through layer $m$ (except for those starting at that layer) have been assigned, so this ensures that the total weight of all chains passing through $L_m$ is $\binom{n}{m}$. Since the weights of all chains starting at one layer are the same, also the induced weights of all points from the same layer are the same, this ensures that the induced weights of all points are $1$.

Now we need only need to prove that all the weights assigned this way are positive, i.e. $W(i) > 0$ for all $i$.

We will do this separately for the outer and inner layers, defined below.

\begin{definition}
    Let us call $L_m$ an \emph{outer layer} if $|m-\frac{n}{2}| \geqslant \frac{k}{2}$. Otherwise, if $|m-\frac{n}{2}| < \frac{k}{2}$, we call $L_m$ an \emph{inner layer}.
\end{definition}

We note that a layer is outer if and only if there are no non-symmetric chains from $\mathcal{C}$ starting on the other side of the cube that pass through it, and inner if there exist such chains. The chains starting at outer layers are all of length $k+1$, and those starting at inner layers are all symmetric and shorter than $k+1$ (for $n$ and $k$ of the same parity, the chains starting at the innermost outer layers $L_{\frac{n-k}{2}}$/$L_{\frac{n+k}{2}}$ are both symmetric and of length $k+1$). On Figure \ref{d1_drawing1}, the three middle layers are inner, and on Figure \ref{d1_drawing2}, the four middle layers are inner.

We note that the weights are symmetric, i.e. $W(m)=W(n-m)$. Hence, it suffices to prove the claim for lower layers.

\subsection{Weights of chains starting in outer layers are positive}

Let us first show that the weights of chains starting at the outer layers are positive. For $L_m$ an outer lower layer, let us note that the chains passing through $L_{m}$ are the same as the chains passing through $L_{m-1}$, except for:
\begin{itemize}
    \item The chains starting at $L_m$ (as they do not pass through $L_{m-1}$).
    \item If $m \geqslant k+1$, the chains starting at $L_{m-k-1}$ (as they reach $L_{m-1}$ but not $L_m$).
\end{itemize}
Hence, we have:
\begin{equation}
    W(m)=\binom{n}{m} - \binom{n}{m-1} + W(m-k-1)
\end{equation}
(where $W(i)=0$ for $i<0$) and we can show by induction on $m$ that $W(m) > 0$ for $m \geqslant 0$, for $L_m$ an outer layer (as the binomial coefficient $\binom{n}{m}$ are increasing for $m < \frac{n}{2}$).

\subsection{Weights of chains starting in inner layers are positive}

Showing that the weights of chains starting at inner layers are positive is a bit more tricky, as these layers have the additional contributions from the other side of the cube.

Let us define:
\begin{equation}
    U(m)= \begin{cases}
        \displaystyle \sum_{\substack{0 \leqslant i \leqslant m, \\ i \equiv m \ (\modd k+1)}} \binom{n}{i} - \binom{n}{i-1}, & \text{if } m\leqslant\frac{n}{2}, \\
         \displaystyle\sum_{\substack{m \leqslant i \leqslant n, \\ i \equiv m \ (\modd k+1)}} \binom{n}{i} - \binom{n}{i+1}, & \text{if } m>\frac{n}{2}.
    \end{cases}
\end{equation}
(the reason for different definitions on both sides of the cube is that $\binom{n}{k}=\binom{n}{n-k}$, so this ensures that $U(m)=U(n-m)$ and $U(m)=W(m)$ for outer layers).

For $L_m$ an outer layer we have $W(m)=U(m)$. We will now show that for $L_m$ an inner layer we have $W(m)=U(m)-U(m+k)$ (where, since $L_{m+k}$ is an outer layer, we have $W(m+k)=U(m+k)$).

Again, it suffices to show this for the lower layers. It is relatively clear for the outermost inner layer, i.e. $L_m$ for $m= \Big\lceil \frac{n-k+1}{2} \Big\rceil$, as the only chains starting in the upper side of the cube passing through this layer are those starting at $L_{m+k}$.

For further inner layers, we can show this by induction: the chains passing through $L_{m+1}$ are exactly the chains passing through $L_{m}$, except for the chains starting at $L_{m+1}$, $L_{m-k}$ and at $L_{m+1+k}$. Hence, we have:
\begin{equation}
    W(m+1)=\binom{n}{m+1}-\binom{n}{m} - W(m+1+k)+W(m-k).
\end{equation}
By the inductive hypothesis, we have $W(m)=U(m)-U(m+k)$ and also, since these layers are outer, $W(m+1+k)=U(m+1+k)$ and $W(m-k)=U(m-k)$. Since, from the definition of $V$, we have $\binom{n}{m+1}-\binom{n}{m} + U(m-k)=U(m+1)$, we get that indeed $W(m+1)=U(m+1)-U(m+1+k)$.

Now we only need to show that for $L_m$ an inner layer, we have $U(m)-U(m+k)>0$. However,
\begin{multline}
    U(m)-U(m+k)\\
    =\sum_{\substack{0 \leqslant i \leqslant m, \\ i \equiv m \ (\modd k+1)}} \bigg( \binom{n}{i} - \binom{n}{i-1}\bigg) - \sum_{\substack{m+k \leqslant i \leqslant n, \\ i \equiv m+k \ (\modd k+1)}} \bigg( \binom{n}{i} - \binom{n}{i+1}\bigg)\\
    =\sum_{\substack{0 \leqslant i \leqslant n, \\ i \equiv m \ (\modd k+1)}} \binom{n}{i} - \sum_{\substack{0 \leqslant i \leqslant n, \\ i \equiv m-1 \ (\modd k+1)}} \binom{n}{i},
\end{multline}
which is positive for $\frac{n-k}{2} < m \leqslant \frac{n}{2}$. It is a well-known fact that can be proven by induction on $n$ (it is presented as Lemma \ref{layers_modulo_largest}). This finishes the proof in the $d=1$ case.

\section{The \texorpdfstring{$d=2$}{d=2} case}
\label{section_d2_overview}
In this section, we prove Theorem \ref{the_main_theorem}.

We note that the special case $k=n$ can be solved in exactly the same way but with slightly different (simpler) details, as all layers behave as outer layers (defined in Definition \ref{definition_outer_d2}). To make the presentation more readable, we will not mention this special case separately in the course of the proof.

\subsection{Definitions and notation}

First, let us start with some definitions and notation, which are natural generalisations of those from Section \ref{section_d1}.

\begin{definition}
    For $x \in \{0, 1, 2\}^n$, let the \emph{type} of $x$ be a tuple $(a, b, c)$ such that $x$ has $a$ coordinates equal $0$, $b$ coordinates equal $1$ and $c$ coordinates equal $2$ (where we must have $a+b+c=n$).

    Later, we often speak about \emph{type $(a, c)$,} meaning type $(a, n-a-c, c)$.
\end{definition}

\begin{notation}
    Let us denote the multinomial coefficients as follows:
    \begin{equation*}
        \binom{n}{a_1, a_2, \dots, a_k} = \frac{n!}{a_1! \cdot \dots \cdot a_k! \cdot (n-a_1 - \dots - a_k)!},
    \end{equation*}
    for $0 \leqslant n, a_1, \dots, a_k, n-a_1-\dots -a_k$.
\end{notation}

We note that there are $\frac{n!}{a!b!c!}=\binom{n}{a,b}$ elements of type $(a, b, c)$ and that all elements of type $(a, b, c)$ lie in layer $b+2c$. We also note that two elements of the same type differ only by a permutation of the coordinates (since they have the same number of $0$s, $1$s and $2$s).

\begin{definition}
    Again, let $L_m$ denote the set of all $x \in \{0, 1, 2\}^n$ with $|x|=m$ for $0 \leqslant m \leqslant 2n$, and we call this set a \emph{layer.}

    Let $L_m$ with $m \leqslant n$ be called \emph{lower layers} and $L_m$ with $m \geqslant n$ be called \emph{upper layers.}

    The layer $L_n$ in the middle of the cube will be called the \emph{middle layer.}
\end{definition}

\begin{notation}
    For $x, y \in \{0, 1, 2\}^n$, let us write $x \leqslant y$ if and only if $x_i \leqslant y_i$ for all coordinates $1 \leqslant i \leqslant n$.

    Moreover, let us write $x \prec_j y$ if $x_i=y_i$ for all $1 \leqslant i \leqslant n, i \neq j$ and $y_j=x_j+1$. We sometimes also write $x \prec y$ if $x \prec_j y$ for some $j$. Again, we note that this relation isn't transitive.
\end{notation}

\begin{definition}
\label{basic_definition}
    We call an ordered tuple $(x_1, \dots, x_l)$ of elements of $\cube$ a \emph{chain} if $x_1 \prec x_2 \prec \dots \prec x_l$.

    We call a chain $(x_1, \dots, x_l)$ with $x_{i} \prec_{j_i} x_{i+1}$ for some $j_1, \dots, j_{l-1}$ \emph{basic} if all the following conditions are met (where $x_{i,j}$ denotes the $j^{\text{th}}$ coordinate of element $x_i$):
    \begin{itemize}
        \item for all $2 \leqslant i \leqslant l-1$ we must have either $x_{i,j_{i-1}}=2, x_{i, j_{i}}=0$ or $j_{i-1}=j_i$,
        \item $x_{1, j_1}=0$,
        \item $x_{l, j_{l-1}}=2$.
    \end{itemize}

    For a chain $x_1 \prec_{j_1} x_2 \prec_{j_2} \dots \prec_{j_{l-1}} x_l$, we call the \emph{width} of that chain the number of different elements of set $\{j_1, j_2, \dots, j_{l-1}\}$.
\end{definition}

For example, the following is a basic chain in $\{0, 1, 2\}^5$ of width $2$: $(0, 1, 1, 0, 0),$ $(0, 1, 1, 1, 0),$ $(0, 1, 1, 2, 0),$ $(1, 1, 1, 2, 0),$ $(2, 1, 1, 2, 0).$

We note that if we have a chain $(x_1, \dots, x_l)$, where $x_1 \in L_m$, then $x_2 \in L_{m+1}, \dots, x_l \in L_{m+l-1}$. Moreover, we note that if $x_i$ is of type $(a, b, c)$, then $x_{i+1}$ is either of type $(a-1, b+1, c)$ or $(a, b-1, c+1)$.

\begin{definition}
    We call a chain $(x_1, \dots, x_l)$ with $x_1 \in L_m$ (so $x_l \in L_{m+l-1}$) \emph{symmetric} if and only if $m=2n-(m+l-1)$.
\end{definition}

We note that a basic chain whose first element is of type $(a, b, c)$ is symmetric if and only if its last element is of type $(c, b, a)$.

\subsection{Plotting the chains in the cube}

To make our proofs and reasoning more intuitive, it is useful to introduce a way of plotting the layers in $\cube$. The types will be represented by unit squares in a coordinate system, where the vertical coordinate represents the type's number of $0$s ($a$) and the horizontal coordinate represents the number of $2$s ($c$). Of course, both of these coordinates start from $0$.

For example, the uppermost square on Figure \ref{figure_types} represents type $(9, 0, 0)$, the two squares below it represent types $(8, 1, 0)$ (the left one) and $(8, 0, 1)$ (the right one), etc.

All types from one layer lie on diagonals going up and right (as they have constant difference $b+2c=n+c-a$) - for example, the types forming the middle layer are marked grey on the picture.

Types along diagonals going down and right have the same value of $b$ (as they have constant $a+c$).

We note that our plot is a projection of $\{a, b, c \geqslant 0: a+b+c=n\}$ from a three-dimensional space $\mathbb{Z}^3$ onto the plane $a,c$.

\begin{figure}
    \centering
    \includegraphics[width=0.4\linewidth]{Types.png}
    \caption{$n=9$}
    \label{figure_types}
\end{figure}

We will later refer to this way of plotting the types as a \emph{staircase diagram.}

\subsection{Weighted chains}

Let $A$ be any set satisfying the problem's conditions (i.e. for any $x \leqslant y$ in $A$ we must have $x_i < y_i$ for more than $k$ different coordinates $i$) and
\begin{equation}
    B = \{x: |x| \equiv n \ (\modd 2k+1) \}.
\end{equation}

We note that $B$ satisfies the problem's condition, and that any set satisfying the condition can have at most one element in any chain of width $k$ or smaller.

Let
\begin{multline}
    \mathcal{C} = \Big\{ \gamma : \gamma \text{ is a basic chain in $\cube$ of width $k$ and length $2k+1$} \Big\} \ \cup\\
    \Big\{ \gamma : \gamma \text{ is a basic, symmetric chain in $\cube$ of width at most $k$ and length less than $2k+1$} \Big\}
\end{multline}
(where \emph{basic} chains were defined in Definition \ref{basic_definition}).
We note that $\mathcal{C}$ contains singleton chains with all elements of the middle layer. We also note that $B$ has exactly one element in all chains in $\mathcal{C}$.

Moreover, we note that each point is an end of at least one chain, whose other end is not closer to the middle (this fact will be important later when assigning weights recursively).

\begin{lemma}
\label{b11111_can_be_0}
    If there exist strictly positive weights $w_{\gamma}$ for all $\gamma \in \mathcal{C}$ (except for $\gamma=((1, 1, \dots, 1))$, which is allowed to have a weight equal to $0$), such that for any $x \in \cube$
    \begin{equation*}
        w_x := \sum_{\substack{\gamma \in \mathcal{C}: \ x \in \gamma}} w_{\gamma} = 1,
    \end{equation*}
    then the set $B$ is the unique maximal set in $\cube$ satisfying the conditions.
\end{lemma}

First we will prove the lemma, and in the rest of the section we will show that such an assignment of weights exists. In the next subsection, we define weights as required in the lemma, and in the further three subsections we show that such weights are indeed non-negative.

\begin{remark}
    We note that, at the first sight, it is by no means obvious that taking only basic chains should be enough to make such an assignment of weights possible (there are many non-basic chains of width $k$ with exactly one element of $B$, for example a chain passing through elements of types: $(a, b, c),$ $(a-1, b+1, c),$ $(a-2, b+2, c),$ $\dots,$ $(a-k, b+k, c),$ $(a-k, b+k-1, c+1),$ $\dots,$ $(a-k, b, c+k)$). We noticed that the basic chains should be enough after making simulations in Python. Since taking as few chains as possible (as long as they allow us an appropriate assignment of weights) appears to make the analysis easier, we decided to just take the basic chains.
\end{remark}

\begin{proof}
    First, we will show that $B$ is the largest possible such set, and then we will show that no other set of equal size satisfies the conditions.

    Assuming that a specified assignment of weights is possible, we can easily finish the proof by observing (precisely as in the case $d=1$) that we have:
    \begin{multline}
        |A| = \sum_{a \in A} 1 = \sum_{a \in A} w_a = \sum_{a \in A} \sum_{\substack{\gamma \in \mathcal{C}: \\ a \in \gamma}} w_{\gamma} = \sum_{\substack{\gamma \in \mathcal{C}: \\ A \cap \gamma \neq \emptyset}} w_{\gamma} \\ \leqslant \sum_{\substack{\gamma \in \mathcal{C}}} w_{\gamma} = \sum_{\substack{\gamma \in \mathcal{C}: \\ B \cap \gamma \neq \emptyset}} w_{\gamma} = \sum_{b \in B} \sum_{\substack{\gamma \in \mathcal{C}: \\ b \in \gamma}} w_{\gamma} = \sum_{b \in B} w_b = |B|.
    \end{multline}
    Now, let us show how the weights of all chains (except for the singleton chain consisting of the point $(1, 1, \dots, 1)$) being positive imply the uniqueness of $B$ as the optimal set. Let us suppose that $|B'|=|B|$ for another set $B'$ satisfying the condition.
    
    Since $B'$ must contain exactly one element from each chain in $\mathcal{C}$ with positive weight, in particular it needs to contain all elements of the middle layer, except possibly $(1, 1, \dots, 1)$ (as $\mathcal{C}$ contains singleton chains consisting of these). However, this means that $B'$ must also contain $(1, 1, \dots, 1)$, as it is maximal and for any element $x \in \cube$ such that $x \leqslant (1, 1, \dots, 1)$ and $x$ with at most $k$ coordinates different from $1$, we also have $x \leqslant y$ for some $y$ from the middle layer of type $(1, n-2, 1)$. Hence, $L_n \subseteq B'$.
    
    Now similarly to the even case for $d=1$, for any $x \in L_{n-2k-1}$, there exists a chain in $\mathcal{C}$ starting at $x$ and finishing at $L_{n-1}$. $B'$ must contain an element of this chain, but it cannot contain any of its elements other than $x$, as this would violate the condition. Hence, $x \in B'$ and so $L_{n-2k-1} \subseteq B'$. We can repeat this reasoning, ultimately finding that $B \subseteq B'$, which, together with $|B'| = |B|$, shows the uniqueness of $B$.
\end{proof}

\subsection{Assigning weights}
\label{assigning_weights}

\begin{definition}
\label{definition_outer_d2}
    We call layers $L_m$ \emph{lower} if $m \leqslant n$ and \emph{upper} if $m \geqslant n$.

    Let us call a layer $L_m$ \emph{outer} if either $m \leqslant n-k$ or $m \geqslant n+k$ and \emph{inner} otherwise. Hence, a layer is outer if and only if there are no chains in $\mathcal{C}$ of length $2k+1$ starting on the other side of the cube that pass through it.

    We call a type outer/inner/lower/upper if the layer it belongs to is such. We note that type $(a, b, c)$ is lower if $a \geqslant c$ and upper if $a \leqslant c$.
\end{definition}

\begin{definition}
    Similarly to earlier, we say that a chain $(x_1, \dots, x_l)$ \emph{starts at $x_1$} if $x_1$ is at least as far from the middle layer as $x_l$, and that it \emph{starts at $x_l$} if $x_l$ is at least as far from the middle layer as $x_1$.

    If the chain is symmetric, we can say that is starts at either at $x_1$ or $x_l$.
\end{definition} 

We note that all chains in $\mathcal{C}$ starting at one type $(a, b, c)$ go through the same types - for example, if $(a, b, c)$ is an outer lower layer, these are: $(a, b, c),$ $(a-1, b+1, c),$ $(a-1, b, c+1),$ $(a-2, b+1, c+1),$ $(a-2, b, c+2),$ $\dots,$ $(a-k, b+1, c+k-1),$ $(a-k, b, c+k).$

\begin{figure}
    \centering
    \includegraphics[width=0.4\linewidth]{Chain.png}
    \caption{$n=9, k=2$}
    \label{figure_chain}
\end{figure}

For example, Figure \ref{figure_chain} shows on a staircase diagram the types through which all chains in $\mathcal{C}$ starting at type $(5, 3, 1)$ pass (for $n=9, k=2$).

If $(a, b, c)$ is an inner layer and $b+2c=n-m$, $m<k$, then all chains starting at $(a, b, c)$ in $\mathcal{C}$ (which must be symmetric) go through the types: $(a, b, c),$ $(a-1, b+1, c),$ $(a-1, b, c+1),$ $(a-2, b+1, c+1),$ $(a-2, b, c+2),$ $\dots,$ $(a-m, b+1, c+m-1),$ $(a-m, b, c+m).$

To all chains in $\mathcal{C}$ starting at the same type we will assign the same weight. Let $W(a, b, c)$ be the total weight of all chains starting at $(a, b, c)$. The assignment of equal weights to all chains starting at the same type ensures that all points from the same type have equal weight. We will sometimes write $W(a, c)$ for $W(a, n-a-c, c)$.

We will assign the weights recursively, starting by setting $W(n, 0, 0)=W(0, 0, n)=1$ and then proceeding from the types furthest from the middle towards the ones closest (for two types in the same layer, it does not matter which one would get the weight assigned first, and both orders will lead to the same weights).

We note that due to starting at the furthest types and moving towards the middle, at the time we need to assign the weight to chains starting at $(a, b, c)$, the weights of all other chains passing through any point of that type are assigned, so we can simply set
\begin{equation}
    W(a, b, c) = \binom{n}{a, b} - \sum_{\substack{\gamma \in \mathcal{C}: \ x \in \gamma \\
    \text{ for some $x \in (a, b, c)$},\\
    \ \gamma \ \text{not starting at $x$}}} w_{\gamma} = \binom{n}{a, b} - \sum_{\substack{(a', b', c'): \text{ chains starting at } \\ \text{$(a', b', c')$ pass through $(a, b, c)$}, \\ (a', b', c') \neq (a, b, c)}} W(a', b', c').
\end{equation}
This ensures that for any type $(a, b, c)$ the total induced weight of points of this type is $\binom{n}{a, b}$, so the induced weight of any point of this type is $1$. We now need to show that the weights set in this way are all non-negative. In the following subsection, we show this for outer layers, then we present a few useful formulae and in the final subsection we show this for inner layers.

\subsection{Weights of chains starting in outer layers are positive}
\label{section_d2_outer}

The result from this subsection is strictly easier to prove than the one from the next subsection about inner layers, and the observations from this section will be crucial for the next section.

\begin{notation}
    Below we often speak about \emph{chains starting at $(a, b, c)$} or $W(a, b, c)$, without checking whether $a, b, c \geqslant 0$, applying a convention that when one of $a, b$ or $c$ is negative, there are no chains starting at $(a, b, c)$ and $W(a, b, c)=0$.
\end{notation}

Due to symmetry ($W(a, c)=W(c, a)$), it suffices to prove the result for the chains starting at lower levels. We recall that through these types pass only chains starting at lower types further from the middle (and no chains starting at upper types). Hence, in the rest of this subsection, we will assume that $(a, c)$ is a lower type. In the next short subsection we will show analogous results for upper types.

Our key observation is that the chains passing through types $(a, c)$ and $(a+1, c-1)$ are identical, except for:
\begin{itemize}
    \item The chains starting at $(a, c)$ (as they do not pass through $(a+1, c-1)$.
    \item The chains starting at $(a+1+k, c-1-k)$ and $(a+1+k, c-k)$ (as they reach $(a+1, c-1)$ but not $(a, c)$).
\end{itemize}
Hence, we have:
\begin{equation}
\label{key_observation}
    W(a, c) = \binom{n}{a, c} - \binom{n}{a+1, c-1} - W(a+1, c) \\
    + W(a+1+k, c-1-k) + W(a+1+k, c-k).
\end{equation}
It might be helpful to visualise the formula above on the staircase diagram. On Figure \ref{key_formula_visualisation}, we can see (for $n=13$, $k=2$) which types contribute, according to Equation \ref{key_observation}, to $W(6, 3)$ (which is the type represented by the lower blue cell). Blue cells mean a contribution of the size of the corresponding type and grey cells represent the contribution of the $W$ of the corresponding type. The cells also contain signs $+/-$, according to the sign of the contribution of that cell.

\begin{figure}
    \centering
    \includegraphics[width=0.5\linewidth]{key_formula_visualisation.png}
    \caption{$n=13, k=2$. Contributions for $W(6, 3)$ according to Equation \ref{key_observation}.}
    \label{key_formula_visualisation}
\end{figure}

Our strategy for proving that $W(a, c)>0$ is to unfold the recursive formula for $W(a, c)$ in an appropriate order, so that ultimately we are able to express $W(a, c)$ as a sum of positive terms. It is motivated by tracking unfolding the formulas on a staircase diagram.  

Our first step is to apply the recursive formula to all terms not involving $k$ (so first to $W(a+1, c)$, then to $W(a+2, c)$ that will appear, etc.). After unfolding all such terms, we get:
\begin{multline}
\label{eq_wac_first_unfolded}
    W(a, c) = \sum_{i \geqslant a} \bigg( \binom{n}{i, c} - \binom{n}{i+1, c-1} + W(i+1+k, c-1-k) + W(i+1+k, c-k) \bigg) \cdot (-1)^{a-i} \\ =
    \sum_{i \geqslant a} \binom{n}{i, c} (-1)^{i-a} - \sum_{i \geqslant a} \binom{n}{i+1, c-1} (-1)^{i-a} \\+ \sum_{i \geqslant a} W(i+1+k, c-1-k) (-1)^{i-a} + \sum_{i \geqslant a} W(i+1+k, c-k) (-1)^{i-a}
    \\ = \sum_{i \geqslant a} \binom{n}{i, c} (-1)^{i-a} - \sum_{i \geqslant a+1} \binom{n}{i, c-1} (-1)^{i-(a+1)}\\ + \sum_{i \geqslant a+1+k} W(i, c-1-k) (-1)^{i-(a+1+k)} + \sum_{i \geqslant a+1+k} W(i, c-k) (-1)^{i-(a+1+k)}.
\end{multline}
We can visualise these contributions on a staircase diagram. Figure \ref{step_1_visualisation} shows the contributions for $W(6, 3)$ for $n=13, k=2$. As on Figure \ref{key_formula_visualisation}, blue cells contribute the size of the corresponding type, grey ones contribute the $W$ of the corresponding type, and the sign in a cell represents the sign of a contribution.

We note that usually, the blue and grey vertical segments would be further from each other. Here they neighbour due to very small $k=2$. 

\begin{figure}
    \centering
    \includegraphics[width=0.5\linewidth]{step_1_visualisation.png}
    \caption{$n=13, k=2$. Contributions to $W(6, 3)$ according to Equation \ref{eq_step_1}.}
    \label{step_1_visualisation}
\end{figure}

Since sums of the form as above will appear very often, let us define notation for them:

\begin{notation}
    Let us denote
    \begin{equation*}
        S(a, c) = \sum_{i \geqslant a} \binom{n}{i, c} (-1)^{i-a}
    \end{equation*}
    and
    \begin{equation*}
        R(a, c) = \sum_{i \geqslant a} W(i, c) (-1)^{i-a}.
    \end{equation*}
\end{notation}

On a staircase diagram, $S(a, c)$ is a vertical segment of blue cells (starting at $(a, c)$) with interchanging $+$s and $-$s. $R(a, c)$ is a similar segment, but with grey cells.

Applying it to the previous formula for $W(a, c)$, we get:
\begin{equation}
\label{eq_step_1}
    W(a, c) = S(a, c) - S(a+1, c-1) + R(a+1+k, c-1-k) + R(a+1+k, c-k).
\end{equation}

Our aim will be to show that $S(a, c) - S(a+1, c-1)>0$ and $R(a, c) \geqslant 0$ for any $a>c$. Let us first focus on $R$.

Let us start by unfolding it further, using Equation \ref{eq_step_1}:
\begin{multline}
    R(a, c) = \sum_{i \geqslant a} W(i, c) (-1)^{i-a}\\ = \sum_{i \geqslant a} \Big( S(i, c) - S(i+1, c-1) + R(i+1+k, c-1-k) + R(i+1+k, c-k) \Big)\\
    = \sum_{i \geqslant a} S(i, c) - \sum_{i \geqslant a+1} S(i, c-1) + \sum_{i \geqslant a+1+k} R(i, c-1-k) + \sum_{i \geqslant a+1+k} R(i, c-k)\\
    = \sum_{i \geqslant a} i \cdot \binom{n}{i, c} (-1)^{i-a} - \sum_{i \geqslant a+1} i \cdot \binom{n}{i, c-1} (-1)^{i-(a+1)} \\+ \sum_{i \geqslant a+1+k} i \cdot W(i, c-1-k) (-1)^{i-(a+1+k)} + \sum_{i \geqslant a+1+k} i \cdot W(i, c-k) (-1)^{i-(a+1+k)}.
\end{multline}
The sums above are very similar to the $R(a, c), S(a, c)$ sums, but have different coefficients.
\begin{notation}
    Let us recall the definition of figurate numbers: $P_0(i)=1_{i>0}$ and for $d>0$:
    \begin{equation*}
        P_d(i)=\sum_{1 \leqslant i \leqslant d} P_{d-1}(i)
    \end{equation*}
    (so, for example, $P_1(i)=i$). They also have a closed form: $P_d(i)=\binom{i+d-1}{d}$.

    Let us now define:
    \begin{equation*}
        S(d, a, c)=\sum_{i \geqslant a} P_d(i) \binom{n}{i, c} (-1)^{i-a}
    \end{equation*}
    and:
    \begin{equation*}
        R(d, a, c) = \sum_{i \geqslant a} P_d(i) W(i, c) (-1)^{i-a}.
    \end{equation*}
    In particular, $S(0, a, c)=S(a, c)$ and $R(0, a, c)=R(a, c)$.
\end{notation}
Hence, according to the above, we have:
\begin{equation}
    R(a, c)=S(1, a, c) - S(1, a+1, c-1) + R(2, a+1+k, c-1-k)+ R(2, a+1+k, c-k).
\end{equation}
We can analogously prove that in general, we have:
\begin{equation}
\label{eq_step_2}
    R(d, a, c)=S(d, a, c) - S(d, a+1, c-1) + R(d+1, a+1+k, c-1-k)+ R(d+1, a+1+k, c-k).
\end{equation}
We can visualise this equation on a staircase diagram as Figures \ref{step_2_visualisation_before} and \ref{step_2_visualisation_after} show. The first of these figures visualises $R(d, 6, 4)$ for $n=16$ and $k=2$ (where numbers in cells denote how many times these cells contribute). The second figure shows, what we get after applying Equation \ref{eq_step_2}.

\begin{figure}[h!]
    \centering
    \begin{minipage}{0.5\textwidth}
        \centering
        \includegraphics[scale=0.3]{step_2_visualisation_before.png}
        \caption{$n=16, k=2: R(d, 6, 4)$}
        \label{step_2_visualisation_before}
    \end{minipage}\hfill
    \begin{minipage}{0.5\textwidth}
        \centering
        \includegraphics[scale=0.3]{step_2_visualisation_after.png}
        \caption{$n=16, k=2:$ $R(d, 6, 4)$ after applying Equation \ref{eq_step_2} (only upper-left part visible).}
        \label{step_2_visualisation_after}
    \end{minipage}
\end{figure}

Combining Equations \ref{eq_step_1} and \ref{eq_step_2}, we get that in order to show that $W(a, c)>0$ for $(a, n-a-c, n)$ a lower outer type, it suffices to show that $S(d, a, c) \geqslant S(d, a+1, c-1)$ with a strict inequality for $d=0$ (as $W(a, c)$ is a sum of such expressions).

To do so, we will first find a closed form for $S(d, a, c)$.

\begin{lemma}
\label{s_closed_form}
    For any $d, a, c \geqslant 0$ we have:
    \begin{equation*}
        S(d, a, c) = \binom{n}{c} \cdot \binom{n-c-d-1}{n-a-c}
    \end{equation*}
\end{lemma}

The proof of the lemma is a very simple induction on $a+d$.

Now we would like to show that $S(d, a, c)-S(d, a+1, c-1)\geqslant 0$ with a strict inequality for $d=0$. This is not true in general. However, we can note that at any time $S(d, a, c) - S(d, a+1, c-1)$ appears as we continue to unfold the recursive definition for $W(a, c)$, the invariant $a \geqslant c+d$ holds.

Indeed, since we start with $d=0$ and $a \geqslant c=c+d$ and whenever we use Equation \ref{eq_step_2} to unfold $R$, the value of $d$ in the next $R$ increases by $1$, and the value of $a$ increases by $1+k>1$, the invariant $a>c+d$ will remain true. Hence, it suffices to check the following, trivial lemma.

\begin{lemma}
\label{claim_s_difference_positive}
    For any $d, a, c \geqslant 0$ with $a \geqslant c + d$, we have
    \begin{equation*}
        S(d, a, c) > S(d, a+1, c-1).
    \end{equation*}
\end{lemma}

\subsection{Recursive equations for upper types}

In this brief subsection, we show (without proofs, as these are completely analogous and are easiest seen on a 'staircase plot') recursive equations for outer upper types, analogous to the ones we found in Subsection \ref{section_d2_outer} for outer lower types. In the rest of this subsection, we let $(a, c)$ be an outer upper type.

The main recursive formula for $W(a, c)$, analogous to Equation \ref{key_observation}, is the following:
\begin{equation}
\label{key_observation_prime}
    W(a, c) = \binom{n}{a, c} - \binom{n}{a-1, c+1} - W(a, c+1) + W(a - 1 - k, c + 1 + k) + W(a - k, c + 1 + k).
\end{equation}
Analogously to $S, R$ (which are sums over vertical segments in a staircase plot), we will define $S', R'$ (sums over horizontal segments).
\begin{notation}
    Let
    \begin{equation*}
        S'(d, a, c) = \sum_{i \geqslant c} P_d(i) \binom{n}{a, i} (-1)^{i-c}
    \end{equation*}
    and
    \begin{equation*}
        R'(d, a, c) = \sum_{i \geqslant c} P_d(i) W(a, i) (-1)^{i-c},
    \end{equation*}
    for $P_d(i)$ figurate numbers.
\end{notation}
Using Equation \ref{key_observation_prime}, we can write $W(a, c)$ as follows (analogously to Equation \ref{eq_step_1}):
\begin{equation}
\label{eq_step_1_prime}
    W(a, c) = S'(0, a, c) - S'(0, a-1, c+1) + R'(0, a-1-k, c+1+k) + R'(0, a-k, c+1+k).
\end{equation}
Analogously to the recursive formula from Equation \ref{eq_step_2}, we have
\begin{equation*}
\label{eq_step_2_prime}
    R'(d, a, c) = S'(d, a, c) - S'(d, a-1, c+1) + R'(d+1, a-1-k, c+1+k) + R'(d+1, a-k, c+1+k).
\end{equation*}
Finally, analogously to Lemma \ref{s_closed_form}, we also have a closed form for $S'$:
\begin{lemma}
\label{s_closed_form_prime}
    We have
    \begin{equation*}
        S'(d, a, c) = \binom{n}{a} \binom{n-a-d-1}{n-a-c}.
    \end{equation*}
\end{lemma}

\subsection{Weights of chains starting at inner layers are positive}
\label{section_d2_inner}

Let us define
\begin{definition}
\label{x_definition}
    Let us define $U(a, c)$ to be $0$ if $a, c$ or $ n-a-c<0$, and otherwise, for $a \geqslant c$, recursively as:
    \begin{equation*}
        U(a, c)=\binom{n}{a, c} - \binom{n}{a+1, c-1} - U(a+1, c) + U(a+1+k, c-1-k) + U(a+1+k, c-k).
    \end{equation*}
    For $a<c$, we define $U(a, c)=U(c, a)$.
\end{definition}
    
Since this definition is exactly the same as the recursive definition for $W(a, c)$ from Equation \ref{key_observation}, we note that for $(a, c)$ an outer type we have $W(a, c)=U(a, c)$.

The following lemma expresses $W(a, c)$ for an inner type in terms of $U$ (the reader should note an analogy with the $d=1$ case).

\begin{lemma}
    For $(a, c)$ an inner lower type, we have
    \begin{equation*}
        W(a, c)=U(a, c)-U(a-k, c+k),
    \end{equation*}
    and for $(a, c)$ an inner upper type ,we analogously have
    \begin{equation*}
        W(a, c)=U(a, c)-U(a+k, c-k).
    \end{equation*}
\end{lemma}

\begin{proof}
    It suffices to show this for lower types.

    We recall that a lower type $(a, c)$ is inner if and only if $a-c<k$, and so the outermost inner types have $a-c=k-1$
    
    For the inner types, besides the chains starting at $(a, c), (a+1+k, c-1-k)$ and $(a+1+k, c-k)$ that go through one of $(a, c), (a+1, c-1)$ but not the other, we also have the chains starting (on the other side of the cube) at $(a-k, c+k)$ and $(a-k+1, c+k)$ (the second one provided that $a-c<k-1$) that pass through $(a, c)$ but not $(a+1, c-1)$. Hence, we have: 
    \begin{multline*}
        W(a, c)=\binom{n}{a, c} - \binom{n}{a+1, c-1} - W(a+1, c)\\ + W(a+1+k, c-1-k) + W(a+1+k, c-k) - W(a-k, c+k) - W(a-k+1, c+k) \cdot 1_{a-c<k-1},
    \end{multline*}
    where we recall that since $(a+1+k, c-1-k), (a+1+k, c-k), (a-k, c+k)$ and $(a-k+1, c+k)$ are outer types, their $W$s and $X$s are equal.
    
    Hence, for the outermost inner types (that is, those where $a-c=k-1$), we indeed have $W(a, c)=U(a, c)-U(a-k, c+k)$.
    For $(a, c)$ an inner type closer to the middle, we have by the inductive hypothesis:
    \begin{multline*}
        W(a, c)
        =\binom{n}{a, c} - \binom{n}{a+1, c-1} - W(a+1, c)\\ + W(a+1+k, c-1-k) + W(a+1+k, c-k) - W(a-k, c+k) - W(a-k+1, c+k)\\
        = \binom{n}{a, c} - \binom{n}{a+1, c-1} - \Big( U(a+1, c) - U(a+1-k, c+k) \Big)\\ + U(a+1+k, c-1-k) + U(a+1+k, c-k) - U(a-k, c+k) - U(a-k+1, c+k)\\
        = U(a, c) + U(a+1-k, c+k) - U(a-k, c+k) - U(a-k+1, c+k)\\ = U(a, c)-U(a-k, c+k),
    \end{multline*}
    which proves the lemma.
\end{proof}

Now we will express $U(a, c)$ (at the moment defined recursively) in a form easier to work with.

\begin{notation}
    Let us denote
    \begin{equation}
        l_C(c) := \frac{C-c}{k+1}.
    \end{equation}
    In the contexts where this notation is used, we always have $c \equiv C \ (\modd k+1)$, so $l_C(c)$ is an integer expressing how many multiples of $(k+1)$ $c$ is away from $C$.
\end{notation}

Using inductive relations developed in the previous section (Equations \ref{eq_step_1} and \ref{eq_step_2}), we can show the following lemma:

\begin{restatable}{lemma}{xminusxweightlemma}
\label{x_minus_x_weight_lemma}
    For $(A, C)$ in one of the inner lower layers, we have
    \begin{multline*}
        U(A, C) - U(A-k, C+k)\\
        =\sum_{\substack{c \equiv C \ (\modd k+1)}} \sum_{j \geqslant 0} \binom{l_C(c)}{j} \binom{n}{c+j} \binom{n-c-j-l_C(c)-1}{B-j}\\ - \sum_{\substack{c \equiv C-1 \ (\modd k+1)}} \sum_{j \geqslant 0} \binom{l_{C-1}(c)}{j} \binom{n}{c+j} \binom{n-c-j-l_{C-1}(c)-1}{B-j},
    \end{multline*}
    where $l_C(c)=\frac{C-c}{k+1}$.
\end{restatable}
The proof of the lemma is mostly technical, hence it is moved to Appendix \ref{appendix_proofs}.

\begin{remark}
    We note that the sum in the lemma above may contain binomial coefficients with a negative upper number, where for $n, k \geqslant 0$ we have:
    \begin{equation}
        \binom{-n}{k} = \frac{-n(-n-1)\cdot \dots \cdot (-n-k+1)}{k!}
    \end{equation}
\end{remark}

Performing some further binomial manipulation, we can show that:

\begin{restatable}{lemma}{xminusxotherformula}
\label{x_minus_x_other_formula}
    We have
    \begin{multline*}
        \sum_{\substack{c \equiv C \ (\modd k+1)}} \sum_{j \geqslant 0} \binom{l_C(c)}{j} \binom{n}{c+j} \binom{n-c-j-l_C(c)-1}{B-j}\\ = \sum_{c \equiv C \ (\modd k+1)} \sum_{0 \leqslant i \leqslant b} \binom{n}{c+i}\binom{l_C(c)+B-i}{l_C(c)}\binom{c+l_C(c)+i}{c+l_C(c)} (-1)^{B-i},
    \end{multline*}
    for $l_C(c)=\frac{C-c}{k+1}$.
\end{restatable}

The proof of this lemma is also mostly technical, and as such has been moved to Appendix \ref{appendix_proofs}.

\begin{notation}
    Let us denote:
    \begin{equation}
    \label{fknbc_def}
        F_k(n, B, C) := \sum_{c \equiv C \ (\modd k+1)} \sum_{0 \leqslant i \leqslant b} \binom{n}{c+i}\binom{l_C(c)+B-i}{l_C(c)}\binom{c+l_C(c)+i}{c+l_C(c)} (-1)^{B-i}
    \end{equation}
    (where, naturally, $F_k(n, B, C)=0$ if $B < 0$ or $C < 0$).
\end{notation}

Hence, from Lemmas \ref{x_minus_x_weight_lemma} and \ref{x_minus_x_other_formula} showing that $U(A, C)-U(A-k, C+k)>0$ is equivalent to showing that
\begin{equation*}
    F_k(n, B, C) - F_k(n, B, C-1)>0 \ \ \text{for} \ \ n-k < B+2C \leqslant n.
\end{equation*}

First, we will prove this in the special cases when either $B=0$ or $C=0$, and then prove it in the general case. These three parts of the proof are presented as three lemmas.

\begin{lemma}
\label{b_is_0}
    For $F_k(n, B, C)$ defined as in Equation \ref{fknbc_def}, we have $F_k(n, 0, C)>0$ for any $n, C$ such that $n-k < 2C \leqslant n$.
\end{lemma}

\begin{proof}
    From the definition of $F_k(n, B, C)$, in particular $F_k(n, 0, C)$ equals
    \begin{equation*}
        \sum_{c \equiv C \ (\modd k+1)} \sum_{0 \leqslant i \leqslant 0} \binom{n}{c+i}\binom{l_C(c)+0-i}{l_C(c)}\binom{c+l_C(c)+i}{c+l_C(c)} (-1)^{0-i}
        = \sum_{c \equiv C \ (\modd k+1)} \binom{n}{c}.
    \end{equation*}
    It is a well-known fact (written with a proof in Lemma \ref{layers_modulo_largest}) that for $0 \leqslant C \leqslant \frac{n}{2}$ this is strictly greater than
    \begin{equation*}
        F_k(n, 0, C-1) = \sum_{c \equiv C-1 \ (\modd k+1)} \binom{n}{c},
    \end{equation*}
    which finishes the proof.
\end{proof}

\begin{lemma}
\label{c_is_0}
    For $F_k(n, B, C)$ defined as in Equation \ref{fknbc_def}, we have $F_k(n, B, 0) \geqslant 0$ for any $n, B$ such that $n-k < B \leqslant n$, with equality if and only if $B=n$.
\end{lemma}

We note that the type with $B=n$ is contains only one point, $(1, \dots, 1)$. We do not mind its weight to be $0$, as explained in the proof of Lemma \ref{b11111_can_be_0}.

\begin{proof}
    To prove this, we go back a couple of steps and recall that $F_k(n, B, C)=U(n-B-C, C)-U(n-B-C-k, C+k)$, so
    \begin{equation}
        F_k(n, B, 0)=U(n-B, 0) - U(n-B-k, k)=U(n-B, 0),
    \end{equation}
    where the second equality holds as for $n-k<B$ we have $U(n-B-k, k)=0$. Hence, it suffices to show that for any $0 \leqslant A \leqslant n$, $U(A, 0) \geqslant 0$, with equality if and only if $A=0$.

    From Definition \ref{x_definition} of $U$ we have $U(n, 0)=1$ and for $0 \leqslant A<n$:
    \begin{equation}
        U(A, 0) = \binom{n}{A} - U(A+1, 0).
    \end{equation}
    From this, it can be easily shown by induction that $U(A, 0)=\binom{n-1}{A-1}$, from which the lemma follows.
\end{proof}

We note that since $F_k(n, B, C)=F_k(n, B, n-B-C)$, the above lemma also shows that $F_k(n, B, C) \geqslant F_k(n, B, C)$ for types $(A, B, C)$ with $A=0$ (with equality if and only if $B=n$).

Before we prove the general case, let us introduce a helpful result on symmetry of $F_k(n, B, C)$.

\begin{restatable}{lemma}{symmetrylemma}
\label{symmetry_lemma}
    We have $F_k(n, B, C) = F_k(n, B, A)$ (where, of course, $A+B+C=n$).
\end{restatable}

The proof of the lemma can be found in Appendix \ref{appendix_proofs}. We are now ready to prove the general case.

\begin{lemma}
    For $B, C \geqslant 0$, $n-k < B+2C \leqslant n$ (that is, type $(n-B-C, B, C)$ is in an inner lower layer), we have $F_k(n, B, C) \geqslant F_k(n, B, C-1)$, with equality if and only if $B=n, C=0$.
\end{lemma}

\begin{proof}
    We will prove this by induction on $n=A+B+C$, with base cases $A=0$, $B=0$ or $C=0$ (proved in Lemmas \ref{b_is_0} and \ref{c_is_0}).
    
    First we note that from the binomials coefficient identity $\binom{n}{k}=\binom{n-1}{k}+\binom{n-1}{k-1}$ we have
    \begin{multline*}
        \binom{n}{c+i}\binom{c+l_C(c)+i}{c+l_C(c)}\\
        =\binom{n-1}{c+i}\binom{c+l_C(c)+i}{c+l_C(c)}+
        \binom{n-1}{c+(i-1)}\binom{c+l_C(c)+(i-1)}{c+l_C(c)}\\+
        \binom{n-1}{i+(c-1)}\binom{l_C(c)+i+(c-1)}{l_C(c)+(c-1)}.
    \end{multline*}
    Since if needed, we can also write $B-i$ as $(B-1)-(i-1)$, we note that we have:
    \begin{equation}
        F_k(n, B, C)=F_k(n-1, B, C)+F_k(n-1, B-1, C)+F_k(n-1, B, C-1).
    \end{equation}

    For $B, C>0$ and $n-k+1 < B+2C < n$, applying the inductive hypothesis is relatively straightforward:
    \begin{itemize}
        \item Since $B, C > 0$ and $(n-1) -k < B+2C \leqslant (n-1)$, by the inductive hypothesis we have $F_k(n-1, B, C)>F_k(n-1, B, C-1)$ (as $C>0$, the inequality is sharp).
        \item Since $B-1 \geqslant 0$, $C>0$ and $(n-1) - k < (B-1)+2C \leqslant (n-1)$, by the inductive hypothesis we have $F_k(n-1, B-1, C)>F_k(n-1, B-1, C-1)$.
        \item Since $B, C-1 \geqslant 0$ and $(n-1) -k < B+2(C-1) \leqslant (n-1)$, by the inductive hypothesis we have $F_k(n-1, B, C-1) \geqslant F_k(n-1, B, C-2)$.
    \end{itemize}
    Adding these three inductive inequalities, we get that:
    \begin{multline*}
        F_k(n, B, C)=F_k(n-1, B, C)+F_k(n-1, B-1, C)+F_k(n-1, B, C-1)\\
        > F_k(n-1, B, C-1)+F_k(n-1, B-1, C-1)+F_k(n-1, B, C-2) = F_k(n, B, C-1),
    \end{multline*}
    which is what we need.

    Now we need to separately consider the special cases when $n, B, C$ do not satisfy the extra requirements that make the induction straightforward. The cases when $B=0$ or $C=0$ are dealt with, respectively, in Lemmas \ref{b_is_0} and \ref{c_is_0}.

    In the case when $B, C > 0$ but $B+2C=n$, by Lemma \ref{symmetry_lemma} we have:
    \begin{equation}
        F_k(n-1, B, C)=F_k(n-1, B, (n-1)-B-C)=F_k(n-1, B, C-1).
    \end{equation}
    Applying the inductive hypothesis as before to get that $F_k(n-1, B-1, C) > F_k(n-1, B-1, C-1)$ and that $F_k(n-1, B, C-1) \geqslant F_k(n-1, B, C-2)$ and adding these three inequalities finishes the inductive proof in this case.

    In the case when $B, C > 0$ but $B+2C=n-k+1$, again by Lemma \ref{symmetry_lemma} we have:
    \begin{equation}
        F_k(n-1, B, C-1) = F_k(n-1, B, (n-1)-B-(C-1)) = F_k(n-1, B, C+k-1)=F_k(n-1, B, C-2).
    \end{equation}
    Using the inductive hypothesis to get that $F_k(n-1, B, C) > F_k(n-1, B, C-1)$ and that $F_k(n-1, B-1, C) > F_k(n-1, B-1, C-1)$ finishes the proof in this case.

    This finishes the proof of the lemma in all cases, and so also the proof of Theorem \ref{the_main_theorem}.
\end{proof}

\section{Conclusions and further research}

\subsection{Conclusions}

In this paper, we have solved Problem \ref{the_main_problem}, i.e. shown that for $d=2$, the unique largest subset $A \subseteq \{0, 1, \dots, d\}^n$ such that there are no distinct elements $x, y \in A$ with $x_i \leqslant y_i$ for all $1 \leqslant i \leqslant n$ and $x_i < y_i$ for at most $k$ coordinates $i$ is
\begin{equation}
    A = \Big\{ x \in \{0, 1, \dots, d\}^n : |x| \equiv n \ (\modd 2k+1) \Big\}.
\end{equation}
The main contribution of this paper is the decomposition of the cube $\cube$ into weighted chains of width at most $k$ that are either of maximal length $dk+1$ or symmetric. A large portion of it is showing that the weights of the chains, when assigned recursively to basic chains (defined in Definition \ref{basic_definition}), are indeed non-negative (and, in fact, positive, except for one singleton chain).

Weighted chain decomposition is preserved under a permutation of coordinates, allowing us to consider only types of points, instead of points themselves.

Moreover, the decomposition is explicit and allows us to shift focus from the construction of such a decomposition to proving its properties (i.e., positive weights) with binomial coefficients.

Solving the case $d=1$ with weighted chains, we discovered a new proof of Sperner's theorem.

\subsection{Methods}

Even in the parts of the proof which are more algebra-heavy, keeping a larger picture in mind and the meaning of the manipulated symbols (helped by, for example, staircase diagrams) was very helpful, if not necessary.

While finding a proof, we also made use of simulations in Python, for example to check that assigning non-negative weights to basic chains (as defined in Definition \ref{basic_definition}) is possible (which also verified our hypothesis for small cases) or that our method of proving that the weights of chains starting in outer layers (Subsection \ref{section_d2_outer}) has a chance of working. It also saved us the time of pursuing some proof ideas that could not work (due to being based on false presumptions).

We note that running fast simulations was possible due to the assignment of weights to chains being a straightforward algorithm.

\subsection{Further research}

An obvious next step after proving this result (which we are currently exploring) is showing that our thesis holds for general $d$. Our conjecture is that for $d>2$, the largest subset $A \subseteq \{0, 1, \dots, d\}^n$ with the required properties is
\begin{equation}
    A = \Big\{ x \in \{0, 1, \dots, d\}^n : |x| \equiv \Big\lfloor \frac{nd}{2} \Big\rfloor \ (\modd dk+1) \Big\}.
\end{equation}
We also conjecture that this set is unique for $2 \mid nd$ and that it is one of two unique sets for $2 \nmid nd$ (where the other one consists of elements $x$ with $x \equiv \Big\lceil \frac{nd}{2} \Big\rceil \ (\modd dk+1)$).

We have proved this conjecture asymptotically, i.e. shown that for any $d, k$ constant and $n \rightarrow \infty$, we have
\begin{equation}
    \frac{|A_{k,d}(n)|}{(d+1)^n} = \frac{1}{dk+1} + o(1).
\end{equation}
The main difficulty in generalising the proof for $d=2$ to the general $d$ is that for larger $d$s, there exist types with no chain of width at most $k$ starting at them and passing through a type from the conjectured optimal set.

A potential way to solve this problem might be to consider a broader family of chains (than just the basic ones), with more choice of weights allowed at each step, so that the types with no chain starting there might get the required weight $1$ from other chains starting earlier.

An interesting additional question to ask would be what makes the family of basic chains sufficient for an assignment of non-negative weights to them being possible? For example, a similar family of `anti-basic chains', i.e. chains that, if starting at type $(A, B, C)$, go through types $(A-1, B+1, C)$, $(A-2, B+2, C)$, $\dots,$ $(A-k, B+k, C)$, $(A-k, B+k-1, C+1)$, $(A-k, B+k-2, C+2)$, $\dots,$ $(A-k, B, C+k)$ is not sufficient, and assigning weight to these chains recursively as in Subsection \ref{assigning_weights} would lead to some negative weights being assigned (which we checked in a simulation).

Another interesting line of investigation would be trying different methods for constructing bounded-width (unweighted) chain decomposition of a cube, or showing that it does not exist.

\appendix

\section{More technical lemmas involving binomial coefficients}
\label{appendix_proofs}

\begin{lemma}
\label{layers_modulo_largest}
    For a given $k \geqslant 2$, let us denote
    \begin{equation*}
        B(n, m) := \sum_{\substack{0 \leqslant i \leqslant n, \\ i \equiv m \ (\modd k+1)}} \binom{n}{i}.
    \end{equation*}
    Then for any $n$, we have $B(n, m)>B(n, m')$ if and only if the closest element of
    \begin{equation*}
        L(n, m) := \{0 \leqslant i \leqslant n: i \equiv m \ (\modd k)\}
    \end{equation*}
    to $\frac{n}{2}$ is closer than the closest element of $L(n, m')$.
\end{lemma}

\begin{proof}
    We will prove the lemma by induction on $n$.
    The thesis is clear for $n \leqslant k$, which is the base case. Now, let us consider larger $n$.

    It suffices to prove that for $m \leqslant \frac{n}{2}$ such that $|\frac{n}{2}-m|\leqslant |m+k-\frac{n}{2}|$ (i.e., $m$ is the element of $L(n, m)$ closest to the middle) and $|\frac{n}{2}-m|<|(m-1+k)-\frac{n}{2}|$ (i.e., the closest element of $L(n, m)$ to the middle is closer than such an element of $L(n, m-1)$), we have $B(n, m)>B(n, m-1)$. We note that these conditions put $m$ in the range $\frac{n-k+1}{2} < m \leqslant \frac{n}{2}$.
    
    Since $\binom{n}{i}=\binom{n-1}{i} + \binom{n-1}{i-1}$, the inequality we wish to prove is equivalent to $B(n-1, m) > B(n-1, m-2)$.
    
    For $m$ in the specified range, the element of $L(n-1, m)$ closest to $\frac{n-1}{2}$ is $m$, and the element of $L(n-1, m-2)$ closest to the middle is either $m-2$ (the first case) or $m-2+k$ (the second case).
    
    In the first case, we have $B(n-1, m) > B(n-1, m-2)$ by the inductive hypothesis, and in the second we need to have $\frac{n-1}{2}-m<m-2+k-\frac{n-1}{2}$, which is equivalent to $n-1-k+2<2m$, which is equivalent to $\frac{n-k+1}{2}<m$, which holds for $m$ from the assumption. Hence, in this case we can also apply the inductive hypothesis and finish the proof.
\end{proof}

\xminusxweightlemma*

\begin{remark}
    The proof of this lemma mostly involves binomial coefficient manipulations, which might seem a bit unmotivated. Plotting contributions from different types on a staircase diagram was very helpful for arriving at it.
\end{remark}

\begin{proof}
    We will show that for $(A, C)$ a lower type, we have
    \begin{multline}
    \label{lower_ac_part}
        U(A, C) = \sum_{\substack{c \equiv C \ (\modd k+1), \\ c \leqslant C}} \sum_{j \geqslant 0} \binom{l_C(c)}{j} \binom{n}{c+j} \binom{n-c-j-l_C(c)-1}{B-j}\\ - \sum_{\substack{c \equiv C-1 \ (\modd k+1), \\ c \leqslant C-1}} \sum_{j \geqslant 0} \binom{l_{C-1}(c)}{j} \binom{n}{c+j} \binom{n-c-j-l_{C-1}(c)-1}{B-j}
    \end{multline}
    and that for $(A, C)$ an upper type, we have
    \begin{multline}
    \label{upper_ac_part}
        U(A, C) = \sum_{\substack{c \equiv C \ (\modd k+1), \\ c \geqslant C}} \sum_{j \geqslant 0} \binom{l_{C}(c)-1}{j} \binom{n}{c+j} \binom{n-c-j-(l_{C}(c)-1)-1}{B-j}\\ - \sum_{\substack{c \equiv C+1 \ (\modd k+1), \\ c \geqslant C+1}} \sum_{j \geqslant 0} \binom{l_{C+1}(c)-1}{j} \binom{n}{c+j} \binom{n-c-j-(l_{C+1}(c)-1)-1}{B-j},
    \end{multline}
    which together (when we substitute $A-k, C+k$ for $A, C$ in the second one) imply the lemma.

    Let us first show Equation \ref{lower_ac_part}. To do so, we will mostly use Equations \ref{eq_step_1} and \ref{eq_step_2}. For easier manipulations, let us denote
    \begin{equation*}
        D(d, A, C) := S(d, A, C) - S(d, A+1, C-1).
    \end{equation*}
    For any $I$, we have:
    \begin{multline*}
        U(A, C) = \sum_{0 \leqslant i \leqslant I} \sum_{0 \leqslant j \leqslant i} \binom{i}{j} D(i, A+i(1+k), C-i(1+k)+j)\\
        + \sum_{0 \leqslant j \leqslant I} \binom{I}{j} \bigg( R\Big(I+1, A+(I+1)(1+k), C-(I+1)(1+k)+j\Big)\\ + R\Big(I+1, A+(I+1)(1+k), C-(I+1)(1+k)+j+1\Big) \bigg)
    \end{multline*}
    (for $I=0$, this is simply Equation \ref{eq_step_1}, and for greater $I$s it follows by induction and by applying \ref{eq_step_2}).

    We note that since for $j > i$ we have $\binom{i}{j}=0$, we can change the inner summation from range $0 \leqslant j \leqslant i$ to range $0 \leqslant j$.
    
    Further noticing that for $I$ sufficiently large, $C- Ik < 0$ and for any $j \leqslant I$ we have
    \begin{multline*}
        R\Big(I+1, A+(I+1)(1+k), C-(I+1)(1+k)+j\Big)\\ = R\Big(I+1, A+(I+1)(1+k), C-(I+1)(1+k)+j+1\Big) = 0,
    \end{multline*}
    we get that
    \begin{equation*}
        U(A, C) = \sum_{i \geqslant 0} \sum_{j \geqslant 0} \binom{i}{j} D(i, A+i(1+k), C-i(1+k)+j).
    \end{equation*}
    After applying Lemma \ref{s_closed_form}, we get that
    \begin{multline*}
        U(A, C) = \sum_{i \geqslant 0} \sum_{j \geqslant 0} \binom{i}{j} \binom{n}{C-i(1+k)+j} \binom{n-C+i(k+1)-j-i-1}{B-j}\\ - \sum_{i \geqslant 0} \sum_{j \geqslant 0} \binom{i}{j} \binom{n}{C-1-i(1+k)+j} \binom{n-C+i(k+1)-j-i}{B-j}.
    \end{multline*}
    Changing from summation over $i$ to summation over $c \equiv C \ (\modd k+1)$ (so $i$ becomes $\frac{C-c}{k+1}=l_C(c)$), we further get that
    \begin{multline*}
        U(A, C) = \sum_{\substack{c \equiv C \ (\modd k+1), \\ c \leqslant C}} \sum_{j \geqslant 0} \binom{l_C(c)}{j} \binom{n}{c+j} \binom{n-c-j-l_C(c)-1}{B-j}\\ - \sum_{\substack{c \equiv C-1 \ (\modd k+1), \\ c \leqslant C-1}} \sum_{j \geqslant 0} \binom{l_{C-1}(c)}{j} \binom{n}{c+j} \binom{n-c-j-l_{C-1}(c)-1}{B-j},
    \end{multline*}
    as desired.

    Now, we need to prove Equation \ref{upper_ac_part}.
    Using Equations \ref{eq_step_1_prime}, \ref{eq_step_2} and Lemma \ref{s_closed_form_prime}, analogously to our considerations above, we get that for $(A, C)$ an upper type we have
    \begin{multline*}
        U(A, C) = \sum_{i \geqslant 0} \sum_{j \geqslant 0} \binom{i}{j} \Big( S'(i, A-i(1+k)+j, C+i(1+k)) - S'(i, A-i(1+k)+j-1, C+i(1+k)+1) \Big)\\
        = \sum_{\substack{a \equiv A \ (\modd k+1), \\ a \leqslant A}} \sum_{j \geqslant 0}  \binom{l_A(a)}{j} \bigg( \binom{n}{a+j}\binom{n-a-j-i-1}{B-j} - \binom{n}{a+j-1}\binom{n-a-j-i}{B-j} \bigg)\\
        = \sum_{\substack{a \equiv A \ (\modd k+1), \\ a \leqslant A}} \sum_{j \geqslant 0}  \binom{l_A(a)}{j}\binom{n}{a+j}\binom{n-a-j-i-1}{B-j}\\ - \sum_{\substack{a \equiv A-1 \ (\modd k+1), \\ a \leqslant A}} \sum_{j \geqslant 0}  \binom{l_A(a)}{j}\binom{n}{a+j}\binom{n-a-j-i-1}{B-j}.
    \end{multline*}
    Hence, it suffices to show that
    \begin{multline}
    \label{what_we_need_lemma}
        \sum_{\substack{a \equiv A \ (\modd k+1), \\ a \leqslant A}} \sum_{i \geqslant 0}  \binom{l_A(a)}{i}\binom{n}{a+i}\binom{n-a-i-l_A(a)-1}{B-i}\\ = \sum_{\substack{c \equiv C \ (\modd k+1), \\ c \geqslant C}} \sum_{j \geqslant 0} \binom{l_{C}(c)-1}{j} \binom{n}{c+j} \binom{n-c-j-(l_{C}(c)-1)-1}{B-j}
    \end{multline}
    (where we changed the summation over $j$ on the left-hand side into a summation over $i$, as this will make our calculations clearer later). 
    First, we note that for $c := n-B-a$, we have
    \begin{multline*}
        \binom{l}{i} \binom{n}{a+i}\binom{n-a-i-l-1}{B-i}=\binom{l}{i} \binom{n}{a+i}\binom{n-a-i-l-1}{n-c-a-i}\\
        = \binom{l}{i} \binom{n}{a+i} \sum_{j \geqslant 0} \binom{-l-1}{j} \binom{n-a-i}{n-c-a-j}\\
        = \binom{n}{a+i} \sum_{j \geqslant 0} (-1)^j \binom{l}{i}\binom{l+j}{j} \binom{n-a-i}{n-c-a-j}
        = \sum_{j \geqslant 0} (-1)^j \binom{l+j}{i, j}\binom{n}{a+i, c+j}.
    \end{multline*}
    Hence, the left-hand side of Equation \ref{what_we_need_lemma} equals
    \begin{equation}
    \label{lhs_first_transformation}
        \sum_{\substack{a \equiv A \ (\modd k+1), \\ a \leqslant A}} \sum_{i \geqslant 0} \sum_{j \geqslant 0}  (-1)^j \binom{l_A(a)+j}{i + j}\binom{i+j}{i}\binom{n}{a+i, c+j}.
    \end{equation}
    Since $c=n-B-a$, we can change the summation from $\sum_{\substack{a \equiv A \ (\modd k+1), \\ a \leqslant A}}$ to $\sum_{\substack{c \equiv C \ (\modd k+1), \\ c \geqslant C}}$, where $l_A(a)=\frac{A-a}{k+1}=\frac{c-C}{k+1}=-l_C(c)$. Hence, the left-hand side of Equation \ref{what_we_need_lemma} equals
    \begin{equation*}
        \sum_{\substack{c \equiv C \ (\modd k+1), \\ c \geqslant C}} \sum_{i \geqslant 0} \sum_{j \geqslant 0} (-1)^j \binom{-l_C(c)+j}{i+j}\binom{i+j}{i}\binom{n}{a+i, c+j}.
    \end{equation*}
    Using the fact that for any $p, q$ positive we have $\binom{p-q}{p}(-1)^p=\binom{q-1}{p}$, we can transform it to
    \begin{equation}
    \label{expression_lhs}
        \sum_{\substack{c \equiv C \ (\modd k+1), \\ c \geqslant C}} \sum_{i \geqslant 0} \sum_{j \geqslant 0} (-1)^i \binom{l_C(c)-1+i}{i+j}\binom{i+j}{i}\binom{n}{a+i, c+j}.
    \end{equation}
    We can transform the right-hand side of Equation \ref{what_we_need_lemma}, completely analogously to the way we transformed the left-hand side into \ref{lhs_first_transformation}. After doing so, we get that the right-hand side also can be expressed as Expression \ref{lhs_first_transformation}, which finishes our proof.
\end{proof}

\symmetrylemma*

\begin{proof}
    We may assume that $(A, C)$ is a lower type (so $(C, A)$ is an upper one). Similarly to earlier, the formal proof of this lemma is written down as binomial symbols manipulation, but it can become intuitive from drawing the formulae on a staircase diagram.
    
    In the same way as in the proof of Lemma \ref{x_minus_x_weight_lemma}, we can show that
    \begin{equation*}
        F_k(n, B, C) = 
        \sum_{\substack{c \equiv C \ (\modd k+1)}} \sum_{i \geqslant 0} \sum_{j \geqslant 0} (-1)^j \binom{l_C(c)+j}{i+j}\binom{i+j}{j}\binom{n}{a+i, c+j}
    \end{equation*}
    and that
    \begin{equation*}
        F_k(n, B, A) = 
        \sum_{\substack{a \equiv A \ (\modd k+1)}} \sum_{i \geqslant 0} \sum_{j \geqslant 0} (-1)^i \binom{(l_A(a)-1)+i}{i+j}\binom{i+j}{i}\binom{n}{a+i, c+j}.
    \end{equation*}
    Since $a+c=n-B=A+C$, $l_C(c)=\frac{C-c}{k+1}=\frac{a-A}{k+1}=-l_A(a)$, we have
    \begin{multline*}
        F_k(n, B, C) = \sum_{\substack{a \equiv A \ (\modd k+1)}} \sum_{i \geqslant 0} \sum_{j \geqslant 0} (-1)^j \binom{-l_A(a)+j}{i+j}\binom{i+j}{j}\binom{n}{a+i, c+j}\\
        = \sum_{\substack{a \equiv A \ (\modd k+1)}} \sum_{i \geqslant 0} \sum_{j \geqslant 0} (-1)^i \binom{(l_A(a)-1)+i}{i+j}\binom{i+j}{i}\binom{n}{a+i, c+j} = F_k(n, B, A),
    \end{multline*}
    (where in the second equality we used $(-1)^p\binom{p-q}{p}=\binom{q-1}{p}$) as desired.
\end{proof}

\xminusxotherformula*

\begin{proof}
    First, we note that we have (where in the second equality we use $(-1)^b\binom{-a}{b}=\binom{a+b-1}{b}$)
    \begin{multline*}
        \binom{n-c-j-l_C(c)-1}{B-j}
        =\sum_{0 \leqslant i \leqslant B-j} \binom{n-c-j}{i} \binom{-l_C(c)-1}{B-j-i}\\
        =\sum_{0 \leqslant i \leqslant B-j} \binom{n-c-j}{i} \binom{l_C(c)+B-j-i}{B-j-i} (-1)^{B-j-i}.
    \end{multline*}
    Hence, we also have
    \begin{multline*}
        \sum_{\substack{c \equiv C \ (\modd k+1)}} \sum_{j \geqslant 0} \binom{l_C(c)}{j} \binom{n}{c+j} \binom{n-c-j-l_C(c)-1}{B-j}\\
        = \sum_{c \equiv C \ (\modd k+1)} \sum_{j \geqslant 0} \sum_{0 \leqslant i \leqslant B-j} \binom{n}{c+j, i} \binom{l_C(c)+B-i-j}{B-i-j, j} (-1)^{B-i-j}.
    \end{multline*}
    The above equals (where in the equation, we change summation over $i$ to summation over $h:=i+j$)
    \begin{multline*}
        \sum_{d \equiv c \ (\modd k+1)} \sum_{j \geqslant 0} \sum_{0 \leqslant i \leqslant B-j} \binom{n}{c+j+i} \binom{c+j+i}{c+j} \binom{l_C(c)+B-(i+j)}{l_C(c)-j, j} (-1)^{B-(i+j)}\\
        = \sum_{d \equiv c \ (\modd k+1)} \sum_{j \geqslant 0} \sum_{j \leqslant h \leqslant B} \binom{n}{c+h} \binom{c+h}{c+j} \binom{l_C(c)+B-h}{l_C(c)-j, j}(-1)^{B-h}\\
        = \sum_{d \equiv c \ (\modd k+1)} \sum_{0 \leqslant h \leqslant B} \binom{n}{c+h} \binom{l_C(c)+B-h}{l_C(c)} (-1)^{B-h} \sum_{j \leqslant h} \binom{c+h}{c+j} \binom{l_C(c)}{l_C(c)-j}.
    \end{multline*}
    Using the well-known fact that $\sum_{i} \binom{r}{a+i}\binom{s}{k-i}=\binom{r+s}{a+k}$, we get that the above equals
    \begin{equation*}
        \sum_{c \equiv C \ (\modd k+1)} \sum_{0 \leqslant h \leqslant B} \binom{n}{c+h} \binom{l_C(c)+B-h}{l_C(c)} \binom{l_C(c)+h+c}{l_C(c)+c}(-1)^{B-h},
    \end{equation*}
    as desired (with a local variable denoted $h$ instead of $i$).
\end{proof}

\section{A new proof of Sperner's theorem}
\label{appendix_sperner}
This proof is a special case of our proof from Section \ref{section_d1} for $k=n$.

Let us decompose $\mathbb{P}([n])$ into weighted symmetric chains, so that for any $X \in \mathbb{P}([n])$ the sum of weights of chains passing through $X$ is $1$ and all weights are positive (we call this sum an \emph{induced weight}).

All chains of the same length will have the same weights, so (if we denote by $c_i$ the number of chains of length $i$) we have:
\begin{itemize}
    \item The chains of length $n+1$ have weight $\frac{1}{c_{n+1}}$
    \item The chains of length $n-1$ have weight $\frac{1}{c_{n-1}}(n-1)$
    \item \dots
    \item The chains of length $n+1-2i$ have weight $\frac{1}{c_{n+1-2i}}\Big( \binom{n}{i} - \binom{n}{i-1} \Big)$.
\end{itemize}
We note that this assignment of weights is positive (as $\binom{n}{i}>\binom{n}{i-1}$ for $i \leqslant \frac{n}{2}$). Moreover, for any set $X \in \mathbb{P}([n])$ the sum of weights passing through it is $1$. Indeed, since the chains are symmetric under a permutation of elements of $[n]$, all sets of the same size have the same induced weight, and the sum of induced weights of all sets of size $i$ is clearly $\binom{n}{i}$.

Let $B = \{X \in \mathbb{P}([n]): |X|= \lfloor \frac{n}{2} \rfloor\}$. We note that $B$ has an element in every symmetric chain.
Hence, for any antichain $A \subseteq \mathbb{P}([n])$ we have (denoting for a chain $C$ its weight by $w_C$):
\begin{equation*}
    |A| = \sum_{X \in A} 1 = \sum_{a \in A} \sum_{\substack{C \ \text{chain}, \\ \ X \in C}} w_{C} \leqslant \sum_{C \ \text{chain}} w_C = \sum_{X \in B} \sum_{\substack{C \ \text{chain}, \\ \ X \in C}} w_{C} = \sum_{X \in B} 1 = |B|,
\end{equation*}
which shows that there is no antichain in $\mathbb{P}([n])$ larger than $B$.

Now we only need to show that $B$ is the unique optimal set (together with $\{X \in \mathbb{P}([n]): |X|=\lceil \frac{n}{2} \rceil\}$).

For $n$ even this is clear, since for any $X$ with $|X|=\frac{n}{2}$, the singleton chain $(X)$ has positive weight, so any optimal set must contain it.

For $n$ odd we note that for any $X \subset Y$ with $X=\lfloor \frac{n}{2} \rfloor$, $Y=\lceil \frac{n}{2} \rceil$ the chain $(X, Y)$ has positive weight, so any optimal set must contain exactly one $X$ or $Y$. From this, we can easily get that an optimal set must contain either $\{X \in \mathbb{P}([n]): |X|=\lfloor \frac{n}{2} \rfloor\}$ or $\{X \in \mathbb{P}([n]): |X|=\lceil \frac{n}{2} \rceil\}$ as a subset, so in fact must be equal to one of these sets.

\begin{thebibliography}{99}

\bibitem{katona} Katona G. O. H., Families of subsets having no subset containing another one with small difference. \textit{Nieuw Archief voor Wiskunde}, XX(3):54-67, 1972.

\bibitem{tour_m_part_p_sperner_families} Aydinian H., Czabarka E., Erdős P. L. and Székely L. A., A tour of M-part L-Sperner families. \textit{Journal of Combinatorial Theory, Series A}, 118(2):702-725, 2011.

\bibitem{combinatorics} Bollobás B., Combinatorics. Set Systems, Hypergraphs, Families of Vectors, and Combinatorial Probability. Cambridge University Press, 1986.

\bibitem{nagys} Nagy D. T. and Nagy K., Chain-dependent Conditions in Extremal Set Theory. arXiv:2201.03663, 2023.

\bibitem{symmetric_chain_decomposition} de Bruijn N. G., van Ebbenhorst Tengbergen C. and Kruyswijk D., On the set of divisors of a number. \textit{Nieuw Archief voor Wiskunde, Series 2}, 23:191-193, 1951.

\bibitem{katona_n2} Katona G. O. H., On a conjecture of Erdős and a stronger form Sperner’s theorem. \textit{Studia Sci. Math. Hung.}, 1:59-63, 1966.

\bibitem{kleitman_n2} Kleitman D. J., On a lemma of Littlewood and Offord on the distribution of certain sums. \textit{Math. Z.}, 90:251-259, 1965.

\bibitem{griggs} Griggs J. R., Saturated chains of subsets and a random walk. \textit{Journal of Combinatorial Theory, Series A}, 47(2):262-283, 1988.

\bibitem{katona_independence_number} Katona G. O. H., A generalization of the independence number. \textit{Discrete Applied Mathematics}, 321:1-3, 2022.

\bibitem{dove_extra} Dove A. P., Generalizations of Sperner's Theorem: Packing Posets, Families Forbidding Posets, and Supersaturation. (Doctoral dissertation), 2013.

\end{thebibliography}

\end{document}
\typeout{get arXiv to do 4 passes: Label(s) may have changed. Rerun}
